\documentclass[11pt]{article}
\usepackage{fontspec}
\usepackage[french]{babel}
\usepackage{amsmath}
\usepackage{amssymb}
\usepackage{amsthm}
\usepackage{mathrsfs}
\usepackage{tikz-cd}
\usepackage{enumitem}

\frenchsetup{AutoSpacePunctuation=false}

\newcommand{\printedpage}[1]{\marginpar{\footnotesize p.~#1}}
\newcommand{\msp}[1]{\marginpar{\footnotesize [#1]}}
\newcommand{\D}{\mathbb{D}}
\newcommand{\Ob}{\mathop{\mathsf{Ob}}}
\newcommand{\Dia}{\textsf{Dia}}
\newcommand{\bs}{\backslash}

\begin{document}

\printedpage{44}
% head: A. Grothendieck : Dérivateurs, Ch. I, Édition des Universités de Montpellier II et Paris VII
le passage de $X$ à $X^\circ$ revient à échanger $p_X$ et $q_X$. La fibre de $\Phi(X)$ en $x \in \Ob X$ est $\simeq (X/x)^\circ$, celle en $y^\circ \in \Ob X^\circ$ est $\simeq y\bs X$
\[
\left\{
\begin{aligned}
&p_X^{-1}(\{x\}) \simeq (X/x)^\circ\\
&q_X^{-1}(\{y^\circ\}) \simeq y\bs X
\end{aligned}
\right.
\]

N.B. : $p_X$ est $W$-lisse \texttt{[plutôt $W$-propre]} à fibres $W$-asphériques, donc dans $W^u$ \texttt{[$W^u =$ ensemble des morphismes qui sont dans $W$ et le restent après tout changement de base]}, $q_X$ est $W$-propre à fibres $W$-asphériques, donc dans $W^u$.

\msp{page 74}

On vérifie
\[
\left\{
\begin{aligned}
&\Phi(X)/x \simeq \Phi(X/x) &\qquad& x \in \Ob X\\
&\Phi(X)/x^\circ \simeq \Phi(x\bs X) &\qquad& x^\circ \in \Ob X^\circ
\end{aligned}
\right.
\]

Je dis que $p_X$, $q_X$ sont des $\D$-\textit{équivalences strictes} et \textit{costrictes.} Par les formules précédentes, cela signifie que les $\Phi(X/x)$, $\Phi(x\bs X)$ sont $\D$-contractiles. Donc que si $Z$ a un objet final, $\Phi(Z)$ est $\D$-contractile. Mais considérons
\[
\begin{tikzcd}
\Phi(Z) \arrow[dr, "q_Z"] & \\
& Z^\circ
\end{tikzcd}
\quad ;
\]
il suffit de prouver que c'est une $\D$-équivalence stricte (puisque $Z^\circ$, ayant un objet initial, est $\D$-contractil), donc que les $z_0\bs\Phi(Z)$ sont $\D$-contractiles. On a par les formules précédentes % [sic] contractil -> contractile, adjectif ailleurs écrit « contractile »
\[
z_0\bs\Phi(Z) \simeq \Phi(z_0\bs Z) = \Phi(T) \quad ,
\]
où cette fois $T$ a un objet initial $\emptyset$ et un objet final $e$. Considérons alors l'objet
\[
\emptyset \longrightarrow e
\]
de $\Phi(T)$, c'est un objet final, car pour tout $x \to y$ dans $\Phi(T)$, il y a une unique flèche
\[
\begin{tikzcd}
x \arrow[r] & y \arrow[d] \\
\emptyset \arrow[u] \arrow[r] & e
\end{tikzcd}
\]
% foot: Version 8 juillet 2001

\printedpage{45}
% head: A. Grothendieck : Dérivateurs, Ch. I, Édition des Universités de Montpellier II et Paris VII
\msp{page 75}

Pour résumer

\textbf{Proposition 1 :} \textit{$p_X$ et $q_X$ sont des} \textsf{hot}\textit{-équivalences strictes,} i.e. \textsf{Hot}\textit{-asphériques (donc aussi des $\D$-équivalences strictes si $\Phi(X)$ est dans le domaine de $\D$) à fibres les $(X/x)^\circ$ et $y\bs X$ respectivement (donc à fibres contractiles).}

\textbf{Corollaire 1.} \textit{Pour que $X$ soit $\D$-contractile, il faut et il suffit que $X^\circ$ le soit.}

\hspace*{\parindent}Plus généralement :

\textbf{Corollaire 2.} \textit{Soit $f : X \to Y$. Pour que $f$ soit une $\D$-équivalence, il faut et il suffit que $f^\circ$ le soit.}

\hspace*{\parindent}On utilisera que le diagramme
\[
\begin{tikzcd}
& \Phi(X) \arrow[dl, "p_X"'] \arrow[dr, "q_X"] & \\
X & & X^\circ
\end{tikzcd}
\]
est fonctoriel en $X$. On a donc
\[
\begin{tikzcd}
& \Phi(X) \arrow[rr, "\Phi(f)"] \arrow[dl, "p_X"'] & & \Phi(Y) \arrow[dl, "p_Y"] \arrow[dd, "q_Y"] \\
X \arrow[rr, "f" near end] \arrow[from=1-2, to=3-2, "q_X"' near end, crossing over] & & Y & \\
& X^\circ \arrow[rr, "f^\circ"'] & & Y^\circ
\end{tikzcd}
\]
les deux carrés commutatifs. Comme $p_X$ et $p_Y$ sont des $\D$-équivalences, $f$ est une $\D$-équivalence si et seulement si $\Phi(f)$ l'est, et de même pour $f^\circ$. Ces raisonnements supposent que si $X$ est dans le domaine de $\D$, $\Phi(X)$ aussi. C'est le cas si \Dia{} est formé des catégories ordonnées finies.
% foot: Version 8 juillet 2001

\printedpage{46}
% head: A. Grothendieck : Dérivateurs, Ch. I, Édition des Universités de Montpellier II et Paris VII
\msp{page 76}

\subsection*{$\D$-fibrations}

\textbf{Proposition 1.} \textit{Soit $f : X \to Y$ dans} \Dia\textit{, conditions équivalentes :}

\hspace*{\parindent}a) \textit{Pour tout $y \in Y$, l'inclusion $X_y \to X/y$ est une $\D$-équivalence, et pour toute flèche $y_1 \to y_0$ de $Y$, $X/y_1 \to X/y_0$ est une $\D$-équivalence.}

\hspace*{\parindent}b) \textit{Pour tout coefficient} constant \textit{$F$ sur $Y$, $f_*f^*F$ est localement constant,}\footnote[33]{Par là on entend que pour tout $u : y \to y'$ dans $Y$, $G_{y'} \to G_y$ est \texttt{[un]} isomorphisme. Sans axiomes supplémentaires sur $\D$, il n'est pas clair que cela implique que les ``restrictions'' de $G$ aux $Y/y$ ou $y\bs Y$, soient isomorphes à des coefficients constants. Voir cependant l'axiome \textbf{Der 3} a$'$), page 97, qui assure que c'est pareil.} \textit{et ``se calcule fibre par fibre''} i.e. \textit{l'homomorphisme canonique}
\[
f_*f^*(F)_y \overset{\sim}{\longrightarrow} H^*_{\D}(X_y, F|X_y)
\]
\texttt{[est un isomorphisme]}.

Ces conditions sont équivalentes, et impliquent les conditions de (b) pour tout $F$ \textit{localement} constant sur $Y$.

La première condition dans (a) signifie que pour tout $y \in Y$, et tout $M \in \Ob \mathcal{A}$,
\[
H^*_{\D}(X/y, M) \longrightarrow H^*_{\D}(X_y, M)
\]
est un isomorphisme. Or cette flèche s'interprète comme
\[
f_*(M_X)_y \longrightarrow H^*_{\D}(X_y, M)
\]
est un isomorphisme, \textit{i.e.} que $f_*(M_X)$ se calcule fibre par fibre. Comme $M_X \simeq f^*(M_Y)$, c'est la deuxième condition (b). La deuxième condition dans (a), pour $y_1 \to y_0$ signifie que pour $M$ dans $\mathcal{A}$,
\[
H^*_{\D}(X/y_0, M) \longrightarrow H^*_{\D}(X/y_1, M)
\]
est \texttt{[un]} isomorphisme, or cette flèche s'interprète comme
\[
f_*(M_X)_{y_0} \longrightarrow f_*(M_X)_{y_1} \quad ,
\]
que ce soit toujours un isomorphisme, c'est la première condition (b). Que (b) soit valable aussi pour $F$ localement constant sur $Y$ est immédiat, car c'est une condition de nature locale par rapport à $Y$.

\msp{page 77}

La proposition duale fait appel aux foncteurs
\[
y_0\bs X \longrightarrow y_1\bs X \qquad \text{et} \qquad X_y \longrightarrow y\bs X \quad ,
\]
et aux faisceaux $f_!f^*M$, pour $M$ constant, ou localement constant sur $Y$.
% foot: Version 8 juillet 2001

\printedpage{47}
% head: A. Grothendieck : Dérivateurs, Ch. I, Édition des Universités de Montpellier II et Paris VII
\textbf{Définition 1.} Si les conditions équivalentes de la proposition 1 sont satisfaites, on dit que $f$ est une $\D$-\textit{fibration cohomologique} (ou $\D$-\textit{fibration gauche}). Si les conditions duales sont satisfaites, on dit que c'est une $\D$-\textit{fibration homologique} (ou $\D$-\textit{fibration droite}).

N.B. Les $\D$-fibrations cohomologiques sont stables par tout changement de base qui est un isomorphisme local. Les $\D$-fibrations homologiques par tout changement de base qui est un co-isomorphisme local. Mais ce ne sont pas des notions stables par \textit{tout} changement de base\footnote[34]{? donner exemple!}. Mais la notion de $\D$-fibration cohomologique est stable par changement de base $\D$-lisse, et celle de $\D$-fibration homologique par changement de base $\D$-propre.

\hspace*{\parindent}Je vais essayer de construire un exemple où le changement de base canule, en posant $Y = \boldsymbol{\Delta}_1$. Se donner $X$ sur $\boldsymbol{\Delta}_1$ revient à se donner une sous-catégorie ouverte $X_0$, et son complémentaire fermé $X_1$
\[
\begin{tikzcd}
& X & \\
X_1 \arrow[ur, hook] \arrow[d, dash, dashed] & & X_0 \arrow[ul, hook'] \arrow[d, dash, dashed] \\
1 & & 0 \arrow[ll]
\end{tikzcd}
\quad .
\]
Dire que $X \to Y = \boldsymbol{\Delta}_1$ est une $\D$-fibration cohomologique signifie simplement
\msp{page 78}
que les restrictions
\[
X_0 \overset{j}{\longrightarrow} X \quad , \qquad X_1 \overset{i}{\longrightarrow} X
\]
sont des $\D$-équivalences. Prenons un changement de base $Y' \to Y = \boldsymbol{\Delta}_1$, \textit{i.e.} un $Y'$ avec $Y'_0, Y'_1$
\[
\begin{tikzcd}
& Y' & \\
Y'_1 \arrow[ur, hook] \arrow[d, dash, dashed] & & Y'_0 \arrow[ul, hook'] \arrow[d, dash, dashed] \\
1 & & 0 \arrow[ll]
\end{tikzcd}
\quad .
\]
On voit que la condition $Y'$ lisse sur $\boldsymbol{\Delta}_1$ signifie que pour tout $y'_1 \in \Ob Y'_1$, la catégorie $Y'_0/y'_1$ des objets $y_0$ de $Y'_0$ au dessus de $y'_1$ (\textit{i.e.} munis de $y_0 \to y'_1$) est contractile. On voit que si $Y'_1 \neq \emptyset$ alors $Y'_0 \neq \emptyset$ -- mais un exemple avec $Y'_0 = \emptyset$, où $Y'_1 = \emptyset$, n'est pas intéressant
% foot: Version 8 juillet 2001
\printedpage{48}
% head: A. Grothendieck : Dérivateurs, Ch. I, Édition des Universités de Montpellier II et Paris VII
(car une projection $I \times J \to I$ est toujours une $\D$-fibration cohomologique). Prenons par exemple pour $Y'$
\[
\begin{tikzcd}
& a \arrow[dl] \\
b & \\
& a' \arrow[ul] \\
\underbrace{\qquad}_{\displaystyle Y'_1} & \underbrace{\qquad}_{\displaystyle Y'_0}
\end{tikzcd}
\quad ,
\]
somme amalgamée de deux copies de $Y = \boldsymbol{\Delta}_1$, donc $X' = X \times_Y Y'$ peut se voir comme somme amalgamée $X \amalg_{X_1} X$. Mais ça semble \textit{pas} être un contre-exemple. Il faudrait que $X_1 \to X'$ ne soit \textit{pas} une $\D$-équivalence.

\msp{page 79}

\textbf{Proposition 2.} \textit{Soit $f : X \to Y$ dans} \Dia. \textit{Conditions équivalentes :}

\hspace*{\parindent}a) \textit{Pour tout $Y'' \overset{p}{\to} Y' \to Y$, d'où}
\[
\begin{tikzcd}
X'' \arrow[r, "q"] \arrow[d] \arrow[dr, phantom, "\text{cart}"] & X' \arrow[r] \arrow[d] \arrow[dr, phantom, "\text{cart}"] & X \arrow[d, "f"] \\
Y'' \arrow[r, "p"'] & Y' \arrow[r] & Y
\end{tikzcd}
\quad ,
\]
\textit{si $p$ est une $\D$-équivalence,} \texttt{[alors]} \textit{$q$ aussi.}\footnote[35]{Il faut supposer $Y', Y''$ \textsf{Hot}-asphériques, cf. page 96, et supposer vérifié l'axiome Der 3 $a'$) de la page 97.}

\hspace*{\parindent}b) \textit{$f$ est une $\D$-fibration cohomologique ``universelle'',} i.e. \textit{est une fibration cohomologique, et le reste après tout changement de base.}

\hspace*{\parindent}b$'$) \textit{Condition duale de} (b) \textit{: $f$ est une $\D$-fibration homologique universelle.}\footnote[36]{b) $\Rightarrow$ a) \underline{canulée}, \textit{cf.} annotation marginale,\footnotemark[37] \textit{cf.} contre-exemple page 89. Mais ça marche pour les dérivateurs $\D_i$ ($i \in \mathbb{N}$) associés aux ``localisateurs fondamentaux'' $W_i$, \textit{i.e.} à la théorie des $i$-types d'homotopie.}\footnotetext[37]{L'implication b) $\Rightarrow$ a) pas prouvée, seulement sous forme plus faible :\par \hspace*{\parindent}a$'$) si $Y''$, $Y'$ sont $\D$-asphériques et 1-connexes, alors $X'' \to X'$ est une $\D$-équivalence. (Il suffit même de le supposer pour $Y'$, $Y''$ ayant des objets finaux (pour impliquer b) ou ayant des objets initiaux (pour impliquer b$'$).)}

\textit{a}) $\Rightarrow$ \textit{b}) : Supposons \textit{a}), soit $Y' \to Y$ changement de base, d'où
\[
\begin{tikzcd}
X' \arrow[r] \arrow[d, "f'"'] & X \arrow[d, "f"] \\
Y' \arrow[r] & Y
\end{tikzcd}
\quad .
\]
% foot: Version 8 juillet 2001

\end{document}
